%% Week 46 -- Tuesday session (flipped classroom): questions from Monday,
%% then Project 3.  Companion to week46slides.tex (Monday lecture on
%% decision trees, bagging and random forests) and to
%% doc/BookML/Projects/Project3.tex (first iteration, 2026-11-07, with the
%% updated LLM rules).  Structure decided 2026-11-07: ~20 min open
%% questions from Monday; then the project text is walked through; then
%% the groups (1) decide which of the three paths interests them most and
%% (2) start drafting a proposal -- other data sets than the suggested ones
%% are welcome.  Aim for the end of the session: a sketch of 3-5 slides.
%% The finished sketch is the weekly exercise of week 47 (the last one,
%% deadline Friday November 20, published already in week 46) and is
%% presented (3-5 min per group) in the final flipped classroom of week 47.
%% No cases/insight frames in this deck: the "cases" are the groups' own
%% proposals.  The environments are kept for consistency with the other
%% Tuesday decks.
%% Build:  pdflatex week46tuesday.tex  (twice)
\documentclass[10pt,aspectratio=169]{beamer}
\usetheme{red_plain}
\usepackage[T1]{fontenc}
\usepackage[utf8]{inputenc}
\usepackage{amsmath,amssymb,bm}
\usepackage{graphicx}
\usepackage{listings}
\usepackage{xcolor}
\usepackage{booktabs}
\graphicspath{{../BookFigures/chapter07_trees_and_ensembles/}{./ch7/}{./}}

%% ---- code listings
\definecolor{codebg}{rgb}{0.97,0.97,0.97}
\lstset{language=Python,basicstyle=\ttfamily\scriptsize,backgroundcolor=\color{codebg},
  keywordstyle=\color{blue!70!black}\bfseries,commentstyle=\color{green!45!black},
  stringstyle=\color{red!60!black},showstringspaces=false,frame=single,
  framerule=0pt,xleftmargin=4pt,xrightmargin=4pt,breaklines=true,columns=fullflexible}

%% ---- discussion frames (kept for consistency, unused in this deck)
\newcounter{insight}
\newenvironment{insight}[1][]{%
  \stepcounter{insight}%
  \begin{frame}{Warm-up \arabic{insight}: #1}%
  \setbeamercolor{block title}{bg=structure.fg,fg=white}%
  \setbeamercolor{block body}{bg=structure.fg!8}%
}{\end{frame}}
\newcommand{\question}[1]{\begin{block}{Question}#1\end{block}}
\newcommand{\answer}[1]{\pause\begin{block}{Answer}#1\end{block}}
\newcommand{\hint}[1]{\pause{\small\emph{Hint: }#1}}

%% ---- case frames (unused)
\newcounter{case}
\newenvironment{case}[1][]{%
  \stepcounter{case}%
  \begin{frame}[fragile,environment=case]{Case \arabic{case}: #1}%
}{\end{frame}}

\setbeamertemplate{footline}[frame number]
\setbeamertemplate{blocks}[rounded][shadow=false]
\setbeamercovered{invisible}

\newcommand{\var}{\mathrm{var}}

\title{Week 46, Tuesday Session}
\subtitle{Flipped classroom: your questions on trees and forests, then
  Project 3 -- read the text, choose a path, and start drafting your
  proposal}
\author{Morten Hjorth-Jensen}
\institute{Department of Physics and Center for Computing in Science Education,\\
  University of Oslo, Norway}
\date{}

\begin{document}

\begin{frame}
\titlepage
\end{frame}

\begin{frame}{How today works}
\footnotesize
\begin{block}{The format (as every Tuesday)}
Mondays are lectures; Tuesdays belong to your questions and to concrete
work in groups.  Today the concrete work is the start of \textbf{Project 3},
the last and freest project of the course.
\end{block}
\begin{enumerate}
  \item \textbf{First $\sim$20 min:} your open questions from Monday's
        lecture -- the impurity measures, the CART split, the Gini and
        entropy arithmetic, our own tree code, Eq.~(7.16), bagging, the
        out-of-bag error, the random forest -- discussed in plenum.
  \item \textbf{$\sim$15 min:} the text of Project 3, walked through: the
        deliverables, the updated rules on LLMs, the three paths, what is
        common to all of them.
  \item \textbf{$\sim$15 min, in your groups:} read the three paths and
        decide which one interests you most -- and why.
  \item \textbf{$\sim$35 min, in your groups:} start drafting a
        \textbf{proposal}: problem, data, methods, plan.  Other data sets
        than the ones we suggest are welcome.  We circulate.
  \item \textbf{Last 5 min:} wrap-up.  \textbf{The aim for today is a sketch of
        3--5 slides} of what you want to do.
\end{enumerate}
\begin{block}{The sketch is the weekly exercise of week 47 -- the last one}
The exercise text (\texttt{exercisesweek47.ipynb}) is published already.  You
present the proposal in the final flipped classroom on Tuesday November 17
and hand in the finished version by \textbf{Friday November 20}.  Details on
the last frame.
\end{block}
\end{frame}

%=============================================================================
\section{Questions from Monday}
%=============================================================================

\begin{frame}{Picking up from Monday}
\footnotesize
\begin{columns}[T]
\begin{column}{0.5\textwidth}
\begin{block}{Monday in one sentence each}
\tiny
A tree partitions the feature space into axis-aligned boxes and predicts a
constant in each, grown top-down and greedily.  A split is scored by the
weighted impurity of its children, Eq.~(7.9) -- MSE for regression, Gini
index or entropy for classification -- and CART keeps the largest gain,
Eq.~(7.10), over every feature and every midpoint; the ten points and the
ride data showed the arithmetic, sixty lines of Python reproduced
\texttt{scikit-learn}.  A full-depth tree has zero training error and large
variance.  $\var(\bar f)=\rho\sigma^2+(1-\rho)\sigma^2/B$ (7.16): bagging
shrinks the second term, the random forest lowers $\rho$ with $m\approx\sqrt p$
features per split (7.18), and the out-of-bag error (7.17) comes for free.
\end{block}

\end{column}
\begin{column}{0.48\textwidth}
\begin{block}{Questions we expect -- ask them}
\scriptsize
\setlength{\itemsep}{0pt}
\begin{itemize}\setlength{\itemsep}{0pt}
  \item Why do the Gini index and the entropy almost always pick the same
        split, and when could they differ?
  \item Why only the midpoints between consecutive values?  What does a
        tree predict outside the range of the training data?
  \item Why does adding trees to a forest never overfit, when adding depth
        to one tree does?
  \item Why $m\approx\sqrt p$ -- and what happens at $m=1$ and $m=p$?
  \item Is the out-of-bag error a substitute for cross-validation?
  \item On the cancer data logistic regression won.  So when are forests
        the right choice?
\end{itemize}
\end{block}
\begin{block}{Everything is in \texttt{week46.ipynb}}
\scriptsize
Impurities, the CART scan, \texttt{build\_id3}, the \texttt{DecisionTree},
\texttt{BaggingClassifier} and \texttt{RandomForestClassifier} classes,
every number on the slides.
\end{block}
\end{column}
\end{columns}
\end{frame}

\begin{frame}{Short answers, in case they are needed}
\scriptsize
\begin{itemize}
  \item \textbf{Gini vs entropy.}  Both are strictly concave in the class
        proportions and nearly proportional for two classes ($S/2$ tracks
        $G$); they rank splits the same way unless two candidates are very
        close.  They differ in \emph{values}, never in the tree, on both of
        Monday's examples.
  \item \textbf{Midpoints.}  The criterion is constant between consecutive
        sorted values, so there is nothing to gain in between; outside the
        training range a tree returns the value of the outermost leaf --
        \textbf{no extrapolation}, and a forest inherits this.
  \item \textbf{Error rate as criterion.}  Piecewise linear: it does not
        reward a split that creates a pure child unless the majority vote
        changes ($(400,400)$ example).  Concave measures do, and a pure node
        never needs splitting again.
  \item \textbf{More trees never overfit.}  Eq.~(7.16): $B$ only enters the
        term $(1-\rho)\sigma^2/B$, which decreases; the bias is untouched.
        Depth increases the variance $\sigma^2$ of each tree instead.
  \item \textbf{$m$.}  $m=p$ is bagging ($\rho$ high); $m=1$ picks the split
        feature at random ($\rho$ low, $\sigma^2$ high).  $\sqrt p$ is a rule
        of thumb; the cancer scan showed $\rho$ from $0.68$ to $0.81$ and a
        flat out-of-bag accuracy -- cross-validate $m$ when it matters.
  \item \textbf{OOB vs CV.}  For bagged ensembles, yes, as an estimate of
        the test error of \emph{that} ensemble; it is not a replacement for
        a held-out test set once you have tuned on it.
  \item \textbf{When forests.}  Tabular data with many mixed features,
        interactions, no good parametric form, little time for feature
        engineering; not for one smooth feature, not for extrapolation, not
        when a linear model is already at the Bayes rate.
\end{itemize}
\end{frame}

%=============================================================================
\section{Project 3: the text}
%=============================================================================

\begin{frame}{Project 3: the frame}
\footnotesize
\begin{columns}[T]
\begin{column}{0.5\textwidth}
\begin{block}{Deliverables and deadline}
As in Projects 1 and 2: a report of about 5000 words and 10--15 figures,
written as a scientific report; the GitHub repository with \texttt{Code},
\texttt{README} and the PDF; delivered as a group in Canvas.
\textbf{Deadline: December 11 (midnight).}
\end{block}
\begin{block}{The updated rules on LLMs -- read them}
The guidelines were changed this week.  For the \textbf{text of the report}
only \emph{contribution level 1} is allowed: spell-checking, grammar, minor
phrasing.  Having an LLM draft, rewrite or restructure a section is
\textbf{not acceptable use}.  For the \textbf{code}, levels 0--4 still apply,
with the inline \texttt{LLM-assisted} tag.  The declaration after the
conclusions is required ($-5$ points if missing) and now states the level
(0 or 1) for every text section and the level for every code file.
\end{block}
\end{column}
\begin{column}{0.48\textwidth}
\begin{block}{What is new: you choose}
Projects 1 and 2 prescribed data and methods.  Project 3 offers
\textbf{three paths}, and you pick one -- or propose your own problem,
after writing to us.
\end{block}
\begin{block}{What is the same for every path}
Own code or libraries (say which); the quality measures of Projects 1--2
where they apply; hyperparameters on training/validation data and the test
set once; seeds and documented settings; a link to the literature; a
critical assessment.  The report structure is the one you know.
\end{block}
\begin{block}{Where the text is}
\texttt{doc/Projects/2026/Project3/} in the course repository: PDF, LaTeX
and a notebook with the same text and two starter cells.
\end{block}
\end{column}
\end{columns}
\end{frame}

\begin{frame}{Path A: Solving differential equations with neural networks}
\scriptsize
\begin{columns}[T]
\begin{column}{0.52\textwidth}
\begin{block}{The problem}
The one-dimensional diffusion equation $u_{xx}=u_t$ on $[0,1]$ with
$u(x,0)=\sin(\pi x)$ and $u(0,t)=u(1,t)=0$ -- or another ordinary or
partial differential equation agreed with us (wave, Poisson, coupled ODEs,
the matrix-eigenvalue ODE).
\end{block}
\begin{block}{The parts}
\begin{enumerate}
  \item[a)] the scheme, the analytical solution, the stability condition;
  \item[b)] the explicit scheme at $\Delta x=1/10,\,1/100$, two times, errors;
  \item[c)] neural networks: a trial function with \textbf{hard-coded}
        conditions, and a \textbf{PINN} with soft conditions; against the
        scheme and the analytical solution;
  \item[d)] nodes, layers, activations, collocation points, loss weights,
        optimizer;
  \item[e)] critical assessment.
\end{enumerate}
\end{block}
\end{column}
\begin{column}{0.46\textwidth}
\begin{block}{Who this is for}
You liked weeks 42--43 (Chapter 9), you know finite differences from
FYS3150/4150 or MAT-MEK4270, and you want a problem with an exact answer to
compare against.
\end{block}
\begin{block}{What you need}
Your Project-2 network, or Keras/PyTorch/JAX; automatic differentiation
for $\partial_x$, $\partial_{xx}$ and $\partial_t$ of the network
\emph{output with respect to its inputs}; patience with training
(PINNs are fussy about the learning rate and the loss weights).
\end{block}
\begin{block}{Typical risk}
Spending all the time on getting one PINN to converge and none on part d).
Get the hard-coded version working first.
\end{block}
\end{column}
\end{columns}
\end{frame}

\begin{frame}{Path B: A classification study with several methods}
\scriptsize
\begin{columns}[T]
\begin{column}{0.52\textwidth}
\begin{block}{The problem}
One multiclass (or rich binary) classification data set that is \emph{not}
MNIST, analyzed with \textbf{at least three} methods of the course, among
them \textbf{at least one neural network} (FFNN or CNN) and \textbf{at
least one tree-based ensemble} (random forest, gradient boosting,
XGBoost); logistic regression as the cheap baseline.
\end{block}
\begin{block}{Default data set: Fashion-MNIST}
$70\,000$ grey-scale $28\times 28$ images of clothing in ten classes, the
MNIST format exactly -- your Project-2 code runs unchanged -- but harder:
roughly $0.84$ (logistic) to $0.91$--$0.94$ (CNNs).  Shirt, T-shirt,
pullover and coat get confused, so the confusion matrix means something.
\end{block}
\begin{block}{Other data}
CIFAR-10; UCI tabular sets (covertype, adult, wine quality, spambase);
sequence data for an RNN; Kaggle; your own field -- write to us first.
\end{block}
\end{column}
\begin{column}{0.46\textwidth}
\begin{block}{The parts}
a) the data in context; b) the theory of your methods; c) implementation,
hyperparameter search with cross-validation, times and parameter counts;
d) results on one test set -- accuracy with standard error, confusion
matrix, per-class precision/recall/$F_1$, one-vs-rest ROC and AUC,
misclassified samples, importances or filters -- against the literature;
e) critical assessment.
\end{block}
\begin{block}{Who this is for}
You want the workflow of an applied study: data, baseline, several methods,
honest comparison, published numbers to beat or to explain.
\end{block}
\begin{block}{Typical risk}
Three methods run once each with default settings and a table of three
accuracies.  The points are in the tuning, the uncertainties and the
analysis of \emph{why} the ranking is what it is.
\end{block}
\end{column}
\end{columns}
\end{frame}

\begin{frame}{Path C: The statistics and mathematics behind the methods}
\scriptsize
\begin{columns}[T]
\begin{column}{0.52\textwidth}
\begin{block}{The idea}
Synthetic data where the \emph{truth is known} -- a function plus noise of
known $\sigma^2$, a Gaussian mixture with known Bayes error
$\Phi(-\Delta/2)$ -- so that bias, variance, the Bayes rate, the correlation
between trees and the bias of an error estimator can all be \textbf{derived
and then measured}.
\end{block}
\begin{block}{The parts}
a) the derivations (bias-variance, OLS/Ridge expectation and covariance,
Bayes error, logistic regression as maximum likelihood with the Fisher
information, Eq.~(7.16), what each resampling estimator estimates);
b) linear models and networks: measured against closed-form bias and
variance, the Ridge path, width/epochs/$L_2$, optional double descent;
c) classifiers against the Bayes error, calibration, class imbalance;
d) ensembles against Eq.~(7.16), $\rho$ against $m$, bias-variance of
tree/bagging/forest; e) CV, LOO, bootstrap, .632 and OOB \emph{as
estimators}: their bias and variance, the one-SE rule, nested CV;
f) critical assessment.
\end{block}
\end{column}
\begin{column}{0.46\textwidth}
\begin{block}{Who this is for}
You enjoyed the derivations more than the data sets, and you want to see
Chapters 2, 3, 5, 7 and 8 confirmed -- or refined -- by experiments you
control completely.  Mostly your own code (projects 1 and 2) plus
\texttt{scikit-learn}.
\end{block}
\begin{block}{What it is not}
Not a literature project and not a big-data project: the data are cheap to
generate, the computing is in the repetitions (many training sets, many
bootstrap samples).
\end{block}
\begin{block}{Typical risk}
Measuring everything and deriving nothing -- or the reverse.  The report
must pair each measurement with the formula it tests, and say where they
disagree and why.
\end{block}
\end{column}
\end{columns}
\end{frame}

\begin{frame}{How we read a Project 3 report}
\scriptsize
\begin{columns}[T]
\begin{column}{0.5\textwidth}
\begin{block}{The same grading as before}
Abstract, introduction, theory and methods, implementation and tests,
results and discussion, conclusions and perspectives, the LLM declaration,
references.  Figures with captions that make a point; tables with
uncertainties; code that we can rerun from the README.
\end{block}
\begin{block}{What the freedom adds}
\begin{itemize}
  \item A clear \textbf{question} on page one: what did you set out to find?
  \item A \textbf{justified choice} of data and methods -- why these, what
        is the baseline, what would a sceptic want to see.
  \item A \textbf{comparison} with something external: analytical solution,
        published benchmark, the formula, the Bayes rate.
  \item A \textbf{critical assessment} that goes beyond ``method X was best''.
\end{itemize}
\end{block}
\end{column}
\begin{column}{0.48\textwidth}
\begin{block}{The seven take-aways in the text}
Method fits problem; a number needs a baseline and an error bar; model
selection is not model assessment; flexibility costs variance; averaging
removes variance, not bias; verification before results; reproducibility.
Use them as the checklist for your conclusions.
\end{block}
\begin{block}{Scope}
About five weeks, next to other courses.  A narrow question answered
properly beats a wide one half done.  The proposal you start today is where
you set the scope -- and where we can tell you if it is too much or too
little.
\end{block}
\end{column}
\end{columns}
\end{frame}

%=============================================================================
\section{Group work}
%=============================================================================

\begin{frame}{Group work, step 1 ($\sim$15 min): which path?}
\footnotesize
\begin{block}{Read the three paths in the project text, then answer in your group}
\begin{enumerate}
  \item Which path do you find most \textbf{interesting} -- and is that the
        same as the one you find easiest?  Say both.
  \item What do you \textbf{already have}?  Project-1 regression and
        bootstrap/CV code, the Project-2 network and its MNIST pipeline,
        today's tree and forest code, finite-difference code from other
        courses, Keras or PyTorch experience.
  \item What would you have to \textbf{learn} for each path (PINNs and
        autodiff of network inputs; gradient boosting and CNN tuning;
        Fisher information and the error-estimator study)?  Is that
        learning the kind you want?
  \item \textbf{Compute and data access}: do you have a laptop that can
        train a CNN on $60\,000$ images, or a GPU in Colab?  Is the data set
        you have in mind downloadable without registration and allowed for
        this use?
  \item Who in the group would do what?
\end{enumerate}
\end{block}
\begin{block}{Decide, and write the decision down in one sentence}
``We choose path \_\_ because \_\_; the main thing we must learn is \_\_.''
If you cannot decide between two paths, note what would settle it and ask
us when we come by.
\end{block}
\end{frame}

\begin{frame}{Group work, step 2 ($\sim$35 min): draft the proposal -- the sketch}
\footnotesize
\begin{block}{The aim for today: 3--5 slides, roughly one per item}
\begin{enumerate}
  \item \textbf{The question and the data.}  One sentence stating what you
        want to find out.  The data: what, where from, how many samples and
        features, what the target is, who has analyzed them before (one or
        two references).  For path A: which equation, which conditions,
        what the analytical or reference solution is.  For path C: which
        synthetic setup, with which known quantities.
  \item \textbf{Methods and baseline.}  Which two or three methods, which
        is the baseline, what is own code and what is a library, which
        quality measures, how the hyperparameters will be chosen.
  \item \textbf{Plan and risks.}  A timeline from now to December 11 with
        the division of labour; the one thing most likely to go wrong and
        what you do if it does; the computing you need.
  \item \textbf{(Optional) The figures you expect.}  Sketch the two or three
        key figures of the report before you have the data: what is on the
        axes, what the curves should look like if you are right.
\end{enumerate}
\end{block}
\begin{block}{Format}
Slides (Beamer, PowerPoint, Keynote, Google Slides -- anything that becomes
a PDF).  Sketchy is fine today; the finished version is next week's
exercise.
\end{block}
\end{frame}

\begin{frame}{If you consider a data set of your own: the checklist}
\footnotesize
\begin{columns}[T]
\begin{column}{0.5\textwidth}
\begin{block}{Practical}
\begin{itemize}
  \item Can it be downloaded today, without registration, in a format you
        can read (csv, npz, images in folders)?  Is the licence compatible
        with coursework?
  \item Size: enough samples to train, validate and test (thousands rather
        than hundreds); small enough to iterate on a laptop.
  \item Labels: clean, documented classes or targets; a known class balance.
  \item A documented or canonical train/test split, or a plan for making
        one.
\end{itemize}
\end{block}
\end{column}
\begin{column}{0.48\textwidth}
\begin{block}{Scientific}
\begin{itemize}
  \item At least one published analysis to compare with -- a benchmark
        table, an article, a Kaggle leaderboard with documented methods.
  \item A question that the methods of the course can answer: tabular
        $\to$ trees and ensembles against logistic regression and a network;
        images $\to$ CNN against the rest; sequences $\to$ RNN; a known
        equation $\to$ path A.
  \item Something you understand: domain knowledge makes the discussion.
\end{itemize}
\end{block}
\begin{block}{Then}
Tell us when we come by, or write to us before the week-47 session.  A
data set that passes this checklist is as welcome as Fashion-MNIST.
\end{block}
\end{column}
\end{columns}
\end{frame}

\begin{frame}{What a sketch can look like (an example for path B)}
\tiny
\begin{columns}[T]
\begin{column}{0.5\textwidth}
\begin{block}{Page 1 -- question and data}
\emph{Can a gradient-boosting model on raw pixels compete with a small
convolutional network on Fashion-MNIST, and at what cost?}  Fashion-MNIST:
$70\,000$ $28\times 28$ grey-scale images, 10 balanced classes,
canonical $60\,000/10\,000$ split; Xiao et al.\ 2017, benchmark table on the
Zalando repository (logistic $\approx 0.84$, GB $\approx 0.88$, CNN
$0.91$--$0.94$).
\end{block}
\begin{block}{Page 2 -- methods and baseline}
Baseline: logistic regression (Project 2, zero hidden layers).
Methods: (1) our Project-2 FFNN, 784--128--10, Adam, $L_2$; (2)
\texttt{HistGradientBoostingClassifier} / XGBoost; (3) Keras CNN of week 44.
Hyperparameters by 5-fold CV on $10\,000$ training images, final models on
all $60\,000$, test set once.  Measures: accuracy $\pm$ SE, confusion matrix,
per-class $F_1$, one-vs-rest ROC/AUC, training time, parameter count.
\end{block}
\end{column}
\begin{column}{0.48\textwidth}
\begin{block}{Page 3 -- plan and risks}
Week 47: data pipeline, baseline, FFNN rerun (A). Week 48: boosting tuning
(B), CNN tuning (A). Week 49: results, figures, misclassified samples (all).
Week 50: writing, LLM declaration, README, freeze on Dec 9.  Risk: CNN
tuning too slow on laptops $\to$ Colab GPU, or a fixed small
architecture with only the learning rate and dropout tuned.
\end{block}
\begin{block}{Page 4 -- expected figures}
Accuracy of all four models against training-set size (log $x$-axis);
confusion matrices side by side with the shirt/T-shirt/pullover/coat
block highlighted; AUC per class; accuracy against training time.
\end{block}
\begin{block}{Why this is a good sketch}
A question, numbers to beat, a baseline, a plan with names and dates, a
risk with a fallback, and figures imagined before the data.  Yours can be
rougher today.
\end{block}
\end{column}
\end{columns}
\end{frame}

\begin{frame}{Prompts for the other two paths while you draft}
\footnotesize
\begin{columns}[T]
\begin{column}{0.49\textwidth}
\begin{block}{Path A}
\begin{itemize}
  \item Diffusion equation as given, or another equation?  If another: do
        you have an analytical or a trusted numerical reference solution?
  \item Which library for the networks, and how will you get $u_{xx}$ and
        $u_t$ of the network -- \texttt{jax.grad} twice, \texttt{torch.autograd},
        \texttt{tf.GradientTape}?
  \item Which architectures and activations go into the part-d) scan, and
        how many collocation points?
  \item What is your measure of error -- against the analytical solution at
        two times, over the domain, as a function of training time?
\end{itemize}
\end{block}
\end{column}
\begin{column}{0.49\textwidth}
\begin{block}{Path C}
\begin{itemize}
  \item Which regression function (Runge 1D, 2D, Franke) and which
        $\sigma$?  Which Gaussian mixture, and will the Bayes boundary be
        linear or curved?
  \item How many training sets and bootstrap samples can you afford per
        experiment?  (Hundreds are needed; plan the loops.)
  \item Which derivations of part a) will you write out in full, and which
        will you cite?
  \item Which optional items (double descent, boosting) will you include,
        if any?
\end{itemize}
\end{block}
\end{column}
\end{columns}
\begin{block}{For everyone}
Which figure of the report is the one you would show to prove your point?
Draw it now, with axes and expected curves.
\end{block}
\end{frame}

%=============================================================================
\section{Wrap-up}
%=============================================================================

\begin{frame}{Wrap-up: the sketch becomes the weekly exercise of week 47}
\scriptsize
\begin{columns}[T]
\begin{column}{0.5\textwidth}
\begin{block}{The last weekly exercise}
Finish the sketch you started today into a \textbf{proposal of 3--5
slides}: question and data, methods and baseline, plan and risks, (optional)
the expected figures.  If you propose a data set of your own, include the
checklist items (source, size, licence, split, literature).  Add the short
LLM declaration as in the earlier weekly exercises.  The exercise text is
already in the week-47 folder; hand in the final version as a PDF in
Canvas, one per group, by \textbf{Friday November 20 (midnight)} -- after
you have presented it and heard our reaction.
\end{block}
\begin{block}{The final flipped classroom, Tuesday November 17}
Each group presents its proposal -- the $3$--$5$ slides as they stand that
day, $3$--$5$ minutes -- we and the other groups react, and you leave with a
plan that we have seen.  Revise it with that feedback before Friday.  This
is the last chance to redirect before the project deadline; use it.
\end{block}
\end{column}
\begin{column}{0.48\textwidth}
\begin{block}{Until then}
\begin{itemize}
  \item Read the full text of your path, including the summary of
        requirements and the take-aways.
  \item Download the data and look at them before you promise anything.
  \item Paths A and B: get the baseline running (explicit scheme; logistic
        regression).  Path C: the data generators and the bias-variance loop.
\end{itemize}
\end{block}
\begin{block}{Monday November 16: boosting}
AdaBoost, gradient boosting, XGBoost, and boosting any loss with automatic
differentiation (Sections 7.11--7.15) -- the method that attacks the bias
that today's ensembles leave alone, and a natural third method for path B.
\end{block}
\end{column}
\end{columns}
\end{frame}

\end{document}
